% GENERATED FILE. Edit canonical lessons, then run npm run generate:books.
% canonical-source-hash: fnv1a64:ffb461285e3228b9
% canonical-lessons: LA-C479-pomarium, LA-C479-olea, LA-C479-prelum, LA-C479-insero, LA-C479-vindemia

\chapter{Orchard, Olive tree, Press, To graft, Grape harvest}
\label{ch:la-orchard-olive-tree-press-to-graft-grape-harvest}
\hlchaptermodality{479}

\begin{chapteropening}
\textbf{By the end of this chapter:} \emph{I can talk about orchards and vineyards: an orchard, an olive tree, the press, grafting trees and the grape harvest.}

Complete the last lesson of chapter 479: I can talk about orchards and vineyards: an orchard, an olive tree, the press, grafting trees and the grape harvest.
\end{chapteropening}

\section[pōmārium, pōmāriī]{pōmārium, pōmāriī — an orchard}

\label{lesson:LA-C479-pomarium}



\begin{quote}
Before the new one: say the Latin for to shear, then the Latin for to milk.
\end{quote}

\subsection*{You'll want to know: pōmārium, pōmāriī}
\textbf{pōmārium, pōmāriī} — \textquotedblleft{}an orchard\textquotedblright{}.

\begin{quote}
In pōmāriō arborēs plēnae pōmōrum sunt.
\end{quote}

\begin{quote}
\emph{In the orchard the trees are full of fruit.}
\end{quote}

Second declension, neuter: \emph{pōmārium, pōmāriī n.}, from \emph{pōmum} 'fruit'. A \emph{pōmārius} sells fruit.

\begin{cousinweb}
\emph{pōmārium} is built on \emph{pōmum} 'fruit' with the suffix \emph{-ārium} 'place for' (Lewis \& Short, under \emph{pōmārius}).
\end{cousinweb}

\begin{grammarlens}[title={What you've built}]
One more word for telling, explaining and describing at B1.
\end{grammarlens}

\subsection*{Your turn}
\begin{itemize}
\raggedright
  \item \emph{Say it:} \emph{pōmārium, pōmāriī}
  \item \emph{Say it:} \emph{pōmārium, pōmāriī}, once more
  \item \emph{Say it:} a sentence of your own with \emph{pōmārium, pōmāriī}
  \item \emph{Recall:} say the Latin for to shear, then the Latin for to milk, then say \emph{pōmārium, pōmāriī} again
\end{itemize}

\begin{tcolorbox}[breakable,colback=teal!4,colframe=teal!35!black,title={Before you move on}]
What does pōmārium, pōmāriī mean? (\textquotedblleft{}An orchard\textquotedblright{}.) Say it once more.
\end{tcolorbox}
\section[olea, oleae]{olea, oleae — an olive tree; an olive}

\label{lesson:LA-C479-olea}



\begin{quote}
Before the new one: say the Latin for to milk, then the Latin for an orchard.
\end{quote}

\subsection*{You'll want to know: olea, oleae}
\textbf{olea, oleae} — \textquotedblleft{}an olive tree; an olive\textquotedblright{}.

\begin{quote}
In colle multae oleae sunt.
\end{quote}

\begin{quote}
\emph{On the hill there are many olive trees.}
\end{quote}

First declension, feminine: \emph{olea, oleae f.} A variant of \emph{olīva}.

\begin{cousinweb}
\emph{olea} was borrowed from Greek \emph{elaia} 'olive, olive tree' (EtymDB), like \emph{olīva} and \emph{oleum} 'oil'.
\end{cousinweb}

\begin{grammarlens}[title={What you've built}]
One more word for telling, explaining and describing at B1.
\end{grammarlens}

\subsection*{Your turn}
\begin{itemize}
\raggedright
  \item \emph{Say it:} \emph{olea, oleae}
  \item \emph{Say it:} \emph{olea, oleae}, once more
  \item \emph{Say it:} a sentence of your own with \emph{olea, oleae}
  \item \emph{Recall:} say the Latin for to milk, then the Latin for an orchard, then say \emph{olea, oleae} again
\end{itemize}

\begin{tcolorbox}[breakable,colback=teal!4,colframe=teal!35!black,title={Before you move on}]
What does olea, oleae mean? (\textquotedblleft{}An olive tree; an olive\textquotedblright{}.) Say it once more.
\end{tcolorbox}
\section[prēlum, prēlī]{prēlum, prēlī — a press (for wine or oil)}

\label{lesson:LA-C479-prelum}



\begin{quote}
Before the new one: say the Latin for an orchard, then the Latin for an olive tree.
\end{quote}

\subsection*{You'll want to know: prēlum, prēlī}
\textbf{prēlum, prēlī} — \textquotedblleft{}a press (for wine or oil)\textquotedblright{}.

\begin{quote}
Servī ūvās in prēlō premunt.
\end{quote}

\begin{quote}
\emph{The slaves press the grapes in the press.}
\end{quote}

Second declension, neuter: \emph{prēlum, prēlī n.}, also \emph{praelum}. The heavy beam that squeezed grapes and olives.

\begin{cousinweb}
\emph{prēlum} is built on \emph{premō} 'to press' (EtymDB; Lewis \& Short), from an older \emph{prem-lom}. English \emph{press} comes from the same verb.
\end{cousinweb}

\begin{grammarlens}[title={What you've built}]
One more word for telling, explaining and describing at B1.
\end{grammarlens}

\subsection*{Your turn}
\begin{itemize}
\raggedright
  \item \emph{Say it:} \emph{prēlum, prēlī}
  \item \emph{Say it:} \emph{prēlum, prēlī}, once more
  \item \emph{Say it:} a sentence of your own with \emph{prēlum, prēlī}
  \item \emph{Recall:} say the Latin for an orchard, then the Latin for an olive tree, then say \emph{prēlum, prēlī} again
\end{itemize}

\begin{tcolorbox}[breakable,colback=teal!4,colframe=teal!35!black,title={Before you move on}]
What does prēlum, prēlī mean? (\textquotedblleft{}A press (for wine or oil)\textquotedblright{}.) Say it once more.
\end{tcolorbox}
\section[īnserō, īnserere]{īnserō, īnserere — to graft, to plant in}

\label{lesson:LA-C479-insero}



\begin{quote}
Before the new one: say the Latin for an olive tree, then the Latin for a press.
\end{quote}

\subsection*{You'll want to know: īnserō, īnserere}
\textbf{īnserō, īnserere} — \textquotedblleft{}to graft, to plant in\textquotedblright{}.

\begin{quote}
Agricola rāmum pirī in mālum īnserit.
\end{quote}

\begin{quote}
\emph{The farmer grafts a pear branch onto an apple tree.}
\end{quote}

\emph{īnserō, īnserere, īnsēvī, īnsitum}, from \emph{serō, serere, sēvī} 'to sow'. A different verb \emph{īnserō, īnserere, īnseruī, īnsertum} 'to put in' gives English \emph{insert}.

\begin{cousinweb}
\emph{īnserō} 'to graft' is built from \emph{in-} 'in' and \emph{serō, serere, sēvī} 'to sow, to plant' (EtymDB; Lewis \& Short). It must be kept apart from \emph{īnserō} 'to put in', from \emph{serō, seruī} 'to join'.
\end{cousinweb}

\begin{grammarlens}[title={What you've built}]
One more word for telling, explaining and describing at B1.
\end{grammarlens}

\subsection*{Your turn}
\begin{itemize}
\raggedright
  \item \emph{Say it:} \emph{īnserō, īnserere}
  \item \emph{Say it:} \emph{īnserō, īnserere}, once more
  \item \emph{Say it:} a sentence of your own with \emph{īnserō, īnserere}
  \item \emph{Recall:} say the Latin for an olive tree, then the Latin for a press, then say \emph{īnserō, īnserere} again
\end{itemize}

\begin{tcolorbox}[breakable,colback=teal!4,colframe=teal!35!black,title={Before you move on}]
What does īnserō, īnserere mean? (\textquotedblleft{}To graft, to plant in\textquotedblright{}.) Say it once more.
\end{tcolorbox}
\section[vindēmia, vindēmiae]{vindēmia, vindēmiae — the grape harvest, the vintage}

\label{lesson:LA-C479-vindemia}



\begin{quote}
Before the new one: say the Latin for a press, then the Latin for to graft.
\end{quote}

\subsection*{You'll want to know: vindēmia, vindēmiae}
\textbf{vindēmia, vindēmiae} — \textquotedblleft{}the grape harvest, the vintage\textquotedblright{}.

\begin{quote}
Autumnō vindēmia est; tōta familia ūvās legit.
\end{quote}

\begin{quote}
\emph{In autumn it is the grape harvest; the whole household gathers grapes.}
\end{quote}

First declension, feminine: \emph{vindēmia, vindēmiae f.} English \emph{vintage} comes from it through French.

\begin{cousinweb}
\emph{vindēmia} is built from \emph{vīnum} 'wine, grapes' and \emph{dēmō} 'to take off, to remove' (EtymDB; Lewis \& Short): taking the grapes off the vine. English \emph{vintage} comes from it through French.
\end{cousinweb}

\begin{grammarlens}[title={What you've built}]
One more word for telling, explaining and describing at B1.
\end{grammarlens}

\subsection*{Your turn}
\begin{itemize}
\raggedright
  \item \emph{Say it:} \emph{vindēmia, vindēmiae}
  \item \emph{Say it:} \emph{vindēmia, vindēmiae}, once more
  \item \emph{Say it:} a sentence of your own with \emph{vindēmia, vindēmiae}
  \item \emph{Recall:} say the Latin for a press, then the Latin for to graft, then say \emph{vindēmia, vindēmiae} again
\end{itemize}

\begin{tcolorbox}[breakable,colback=teal!4,colframe=teal!35!black,title={Before you move on}]
What does vindēmia, vindēmiae mean? (\textquotedblleft{}The grape harvest, the vintage\textquotedblright{}.) Say it once more.
\end{tcolorbox}
